% Argumentative expansion. Insert after the introduction of excepcional.tex,
% before its first subsection. Preserves the existing development in full.
\paragraph{Articulation of regional coordinates and closure incidence.}
\label{inc-ampl-articulacion}
The role of this prolongation is understood through a distinction
between three operations: preserving a record, evaluating a coordinate, and
extracting an incidence relation. All three act on information
produced by APP--TRIT--TPK. Their articulation makes explicit the article's
structural thesis: a numerical realization has an operator
provenance, and certain data from that provenance additionally admit a
combinatorial and lattice realization. The equations allow both
paths to be followed with their domains and the information remaining in each.

The starting point is the terminal state with its compatible
prolongations. It preserves the residues and quotients of the two APP
sheets, the TRIT regime and orientation, and TPK transport with its memory.
The closure, propagation, and self-scale regions are determined in that
domain before their readers establish the coordinates
$\pi,e,\varphi$. At the same time, the oriented record of the twelve
windows determines the signed vector $U$. The subsequent appearance of $K$
arises from the integral transformation proved in
\S\ref{subsec:k-hadamard}:
\begin{equation}
 U=\mathcal H_{12}K,\qquad
 K=\frac14\mathcal H_{12}U,\qquad
 \mathcal H_{12}^{2}=4I_{12}.
 \label{inc-ampl-registro}
\end{equation}
The congruence defining $\mathcal L_H$ ensures integrality of
the inverse. Thus $K$ is a reversible coordinatization
of the signed record in the relevant sublattice. Its positions preserve
phase order; its differences $D_3K,D_4K$ preserve transverse
relations; the sum $Q(K)$ determines the uniform component. That
structure explains why the record enters both numerical composition
and the selection of a boundary configuration.

The periodic evaluation
\[
 K\longmapsto\kappa_{\mathrm{per}}
 =\frac{\sum_{m=1}^{12}K_m1000^{12-m}}{1000^{12}-1}
\]
keeps the twelve blocks recoverable when base, length, and phase origin
are fixed. Extracting a support performs a different
operation: it preserves membership of each position in a class defined
by a threshold and groups records with the same support. The
distinction between these two maps is decisive. The scalar
precision of $\kappa_{\mathrm{per}}$ and the incidence of $B_K$ are two
specific mathematical contents, related through the complete
record from which they arise.

\paragraph{Alpha normalization and its signed configuration.}
\label{inc-ampl-clausura}
The coordinate of $\alpha$ is determined in the positional route through
composition of the regional blocks and the dodecaphase record. Before
positional division, the oriented combination has components
$s_j=P_j+E_j-\Phi_j-K_j^{\mathrm{per}}$. Normalization satisfies
\begin{equation}
 A_j=s_j+c_{j+1}-1000c_j,\qquad 0\leq A_j<1000.
 \label{inc-ampl-normalizacion}
\end{equation}
This identity simultaneously preserves the normalized block and the
contribution of the transported quotient. The limit of compatible
prefixes determines the scalar coordinate of closure. The
signed configuration, in turn, preserves which positions of the
combination before normalization have a negative defect.
In the initial window, that configuration is the vector $s_0$ of
\eqref{exc:precarry}, and its negative support is $H^-_\alpha$.

Two contents of the same operation are thereby defined. Positional
normalization determines a real coordinate through its compatible
remainders and quotients; the negative support determines a
configuration of positions in the twelve-component cycle. The sign
refers to the balance of records before normalization. Its incidence
reading remains available because the state preserves that balance
together with the normalized result. This articulation corresponds to the
proved positional route; comparison with the analytic vacancy
operators preserves the precise scope established in the chapter on
$\alpha$.

\paragraph{From ternary coding to the marked flag.}
\label{inc-ampl-bandera}
The regions also supply the codes $w_\pi,w_e,w_\varphi$, and
$w_{\rm APP}$ in the inherited frame. The map
\[
 \gamma_{A_W}:w\longmapsto(w,wA_W)
\]
preserves the six-component word as its first coordinate and adds
its transformation by $A_W$. The relation $A_W^2=-I_6$ fixes the
quadratic compatibility of this realization. The next step,
$c\mapsto\operatorname{supp}(c)$, preserves the nonzero positions.
For weight-six words, the support identifies exactly the pair
$\{c,-c\}$, as proved by the minimum distance of the code. Thus,
the sign quotient produces the hexads, whereas the complete code
remains available for operations requiring its symbols.

In this chart, the two constructions of the tetrad coincide:
\begin{equation}
 \{i:K_i\geq729\}
 =H_e\cap P_\pi=B_K=\{4,6,9,10\}.
 \label{inc-ampl-cuaterna}
\end{equation}
The left-hand side arises from an integer comparison of the components
of $K$; the right-hand side, from the intersection of the coded supports of
propagation and closure. The threshold belongs to the declared block
chart and remains explicit in the reader's definition. The
complete record preserves the values on which it acts, even when
different thresholds produce the same subset. The equality therefore
expresses compatibility between specified maps.

The support $H^-_\alpha$ completes that tetrad to a hexad. The
condition of membership in support $P_\pi$, together with the intersection
sizes $(4,3,2)$ relative to $H_e,H_\varphi,H_{\rm APP}$,
selects a unique hexad. The following proofs bring together the two
determinations: that obtained through the sign of $s_0$ and that obtained
through incidence. The result is organized as
\begin{equation}
 \widetilde{\mathcal I}_*
 =\bigl(B_K\subset H^-_\alpha;
         P_\pi,H_e,H_\varphi,H_{\rm APP}\bigr).
 \label{inc-ampl-marco}
\end{equation}
The inclusion is the flag; the four additional supports constitute
its marked frame. Simultaneous permutations transport the frame and
preserve its inclusions and intersections. The relevant information resides
in this transported configuration and its isomorphism class.

The five regions of $\pi$ recover their internal differentiation here.
They share $w_\pi$ as a visible ternary code but preserve distinct emissions,
signatures, and multiplicities. In the binary realization on a
marked dodecad, the transverse region corresponds to the octad whose
intersection with the dodecad is the tetrad; the other four regions
correspond to the lifts of the four hexads of its star. The
negative region is distinguished by $H^-_\alpha$, and chart order
individualizes the three positive regions. This correspondence preserves
the regional roles and distribution $432+4\cdot144=1008$, while
also realizing the incidence partition $1+4$.

\paragraph{From the incidence origin to lattice modification.}
\label{inc-ampl-reticulo}
The marked flag determines a position through an equivariant
set-theoretic operation:
\[
 o_*=(H_e\cap H_\varphi)
       \setminus(H^-_\alpha\cup H_{\rm APP})=\{1\}.
\]
Its role is to distinguish one component among the twelve
entering the lattice realization. The complete ternary code
acts as a glue code in the discriminant module
$(A_2^*/A_2)^{12}$. Its preimage defines $N(C_W)$, and the
self-duality, weights, and distance of the code translate,
respectively, into integrality and determinant properties, parity,
and norm bounds for the glue classes. Here the symbolic information
of the code performs a role that the isolated support no longer performs.

The $A_2$ metric and positive radial chamber are fixed in the lattice
development. Within that chamber, the congruence needed for an
even index-three neighbor has twelve realizations of minimum norm $54$,
permuted by changing the distinguished component. The origin $o_*$
selects
\[
 v_\alpha=4\rho_1+\rho_2+\cdots+\rho_{12},\qquad
 v_\alpha^2=2(4^2+11)=54.
\]
The coefficient $4$ belongs to that solution of the congruence in the
declared chamber. The subscript $\alpha$ records the incidence
provenance of the selection: its content is the frame that entered
closure and now individualizes a lattice direction.

Pairing with $v_\alpha$ determines the index-three sublattice
$N_{\alpha,0}$, and adjoining the class
$v_\alpha/3$ produces the neighbor $N_\alpha$ of
\eqref{exc:vecino}. This modification preserves rank and recovers
unimodularity and evenness; congruence selection eliminates the preceding
roots, and local bounds exclude roots in the new classes.
Identification with Leech appears after those verifications. The
binary realization of the five regions and this lattice construction
are two prolongations of the common frame: the latter directly receives
the ternary code as glue data.

\paragraph{Grading, unit, and automorphisms.}
\label{inc-ampl-automorfismos}
The Frenkel--Lepowsky--Meurman construction is applied to the lattice
$N_\alpha$ once its properties have been verified. It incorporates the oscillator
space, central extension of the lattice, vertex product,
and sector twisted by the lift of negation. The result is
the algebra $V^\natural_{\rm HMT}$ of \eqref{exc:flm}, whose automorphism
group is isomorphic to the Monster. Its action domain is thus
fixed by an algebraic structure with grading and product. Borcherds'
Moonshine theorem subsequently determines the graded traces
with insertions of those automorphisms \cite{flm,borcherds}.

At this level, the unit acquires a precise realization: the
weight-zero component is the vacuum line and contributes the coefficient $1$ of
$q^{-1}$ in the graded character. At the cylinder level,
refinement maps preserve the constant function through
precomposition. Every conservation statement therefore has its
object and its map. The development of the proofs makes this
continuity of construction visible, from arithmetic operations and
memory to incidence, metric, and vertex products.

This articulation places $K$ and the closure of $\alpha$ in the main
argument. $K$ reversibly preserves an integral coordinate of the
history; normalization of $\alpha$ composes the regional records
with that coordinate; the signed configuration determines a flag; and
the marked flag selects the lattice modification on which
Moonshine is realized. The following proofs develop these
maps, distinguishing reversible equivalences, information
quotients, and realizations that add metric or algebraic
structure. A provable relation between structure,
equation, and value is thus expressed, preserving the attribution of each construction
and the scope of its theorems.
